\chapter{Regev optimizations: Fibonacci exponentiation}

The first optimization we are going to talk about has to do with transforming the base in which we represent the numbers in a base that is perhaps more suited when dealing with exponents. To do this, we will entice the amateur mathematics enthusiast by bringing the Fibonacci numbers and Fibonacci sequence into the picture.
To explore this concept, first we will need to talk about Zeckendorf's theorem. The statement is the following:\\

$\textbf{Lemma 15.1 (Zeckendorf)}:$\\
Every positive integer $k$ can be uniquely represented as a sum of non-consecutive Fibonacci numbers.In other word, for every $k \in \mathbb{Z}$:
\begin{equation}
    k = \sum_j c_jFi, \: c_j \in \set{0,1}
\end{equation}

No two consecutive $c_i$ can be 1, because consecutive Fibonacci numbers can be replaced with the next Fibonacci number according to the recursive formula $F_n = F_{n-1} + {F_n-2}$. This means that representing exponents as bits using a Fibonacci base is perfectly valid. We will now describe a way of computing $a^{k}$ using a Fibonacci base. First, we need to compute the Fibonacci representation of k. To do that in a quantum circuit, we use the following greedy algorithm.

$\textbf{Step 1.}$ Compute all Fibonacci numbers up to $k$ using the recursive formula and store it in an array F, starting from $F[0] = 1,F[1] = 2$.\\
$\textbf{Step 2.}$ Initialize an array $F_b$ of 0s with size $m-1$, where $F_m$ is the largest Fibonacci number smaller than $k$\\
$\textbf{Step 3.}$ Let $j \gets m - 1$.\\
$\textbf{Step 4.}$ $F[j] \gets 1$\\
$\textbf{Step 5.}$ $k \gets k - F_m$\\
$\textbf{Step 6.}$ Run binary search on $F$ to compute the largest Fibonacci number smaller than k\\
$\textbf{Step 7.}$ $F_m \gets$ the value computed in Step 6\\
$\textbf{Step 8.}$ $j \gets$ the index of that value in $F$ \\
$\textbf{Step 9.}$Repeat from Step 2 until $k = 0$.\\


F contains the binary representation of k. Now, computing $a^k$ is very easy, because we can use the same logic as modular exponentiation.\\

Assume the greedy algorithm in the previous step leaves us with a Fibonacci representation $k_{\phi}$ consisting of all the $c_j$, and $m$ Fibonacci numbers computed, starting from $F[0] = F_2 = 1$. Let us present the algorithm for Fibonacci exponentiation, as described by Kaliski \cite{Kaliski2017}\\

\begin{algorithm}[H]
\caption{Fibonacci exponentiation}
\textbf{Input:} $a, k, N$, the sequence $c_1,c_2 \dots c_m$\\
\textbf{Output:} $a^k \pmod N$
\begin{algorithmic}[1]
\State $(c,d,e) \gets (1,a,1)$
\For{$j = 1$ to $m$}
    \State $(c,d) \gets (d,cd \pmod N)$
    \If{$c_j == 1$}
        \State $e \gets de \pmod N$
    \EndIf
\EndFor
\end{algorithmic}
\end{algorithm}


Consider the original Regev product $\prod_{i=1}^d b_i^{e_i}$ (the algorithm is also compatible with the $D/2$ shift). Suppose we have a Fibonacci representation of $e_i$ 
\begin{equation}
    e_i = \prod_{j=1}^m c_{i,j} F_j
\end{equation}
where $c_{i,j}$ denotes each of the earlier $c_j$ but now for every base. Then we can plug \textbf{Eq 15.2} into the Regev product and write the resulting expression - which is \textbf{equivalent} to the Regev product, as follows: $$\prod_{j=1}^m(\prod_{i=1}^d b_i^{c_{i,j}})^ {F_j}$$
Denote the inner product as $C_j$. This leaves us with
\begin{equation}
    \prod_{j=1}^m C_j^ {F_j}
\end{equation}

What we need to do here essentially is compute this over all possible $C_j$ in superposition,  raise $C_j$ to the power of the next Fibonacci number, and do this over and over until $F_j > 2^d$. To compute the $C_j$, we can basically do exactly what we did in Chapter 13 for modular multi-exponentiation - precompute the inverses, build the tree, multiply into the accumulator, uncompute the tree. 

Where this logic shines however, is with regards to squaring. In our implementation, we had to use a different algorithm for squaring. However, we don't actually need to square when representing numbers in the Fibonacci base. Consider what happens if we sequentially applied our MOD-PROD gate to some arbitrary  $\ket{a} \ket{a}$: \\

1.$\ket{a} \ket{a^{-1}} \ket{a} \ket{a^{-1}} \to \ket{a} \ket{a^{-1}} \ket{a^2} \ket{(a^{2})^{-1}}$\\
2.$\ket{a} \ket{a^{-1}} \ket{a^2} \ket{(a^{2})^{-1}} \to \ket{a^3} \ket{(a^{3})^{-1}} \ket{a^2} \ket{(a^{2})^{-1}}$\\
3.$\ket{a^3} \ket{(a^{3})^{-1}} \ket{a^2} \ket{(a^{2})^{-1}} \to \ket{a^3} \ket{(a^{3})^{-1}} \ket{a^5} \ket{(a^{5})^{-1}}$\\

Notice the Fibonacci sequence appearing in the exponents! Now we have a reversible base accessing operation that doesn't leave us with dirty qubits in every iteration. More importantly, it seems that this can actually decrease the qubit count using the optimized multi-exponentiation idea from Chapter 7 with the reduced register sizes, because we won't need the extra $O(n^{3/2})$ qubits to perform squaring anymore.

Let us, finally, describe Algorithm 5.2 from Ragavan and Vaikuntanathan\cite{RagavanVaikuntanathan2023} in terms of what we wrote earlier. It is important to note that this is applied to the shifted exponent $e_i = c_i + D/2$.

\begin{algorithm}[H]
\caption{Fibonacci multi-exponentiation (top-down)}
\textbf{Input:} Initial quantum state $\ket{e}$\\
\textbf{Output:} State consisting of a superposition $\prod_{i=1}^d b_i^{e_i} \pmod N$
\begin{algorithmic}[1]
\State Compute $c_{i,j}$ in superposition according to our initial algorithm.
\State $R \gets 1$
\For {$j = 0$ to $m$}
    \State Compute $\ket{C_j}$ in the manner we described in Chapter 13, entangling $the e_i$ in a superposition.
    \State Run the multiplication circuit on $C_j$ and $C_j^{-1}$ j times to get the next Fibonacci base
    \State $R \gets RC_j^{F_j}$
    \State Run the multiplication circuit again $j$ times to turn $\ket{C_j^ {F_j}} \ket{C_j^{-F_j}} \to \ket{C_j} \ket{C_j^{-1}}$  
    \State Uncompute the intermediate nodes used to compute $\ket{C_j}$ using $\ket{C_j^{-1}}$
\EndFor
\State $R_1$ is now entangled with $\ket{e}$, and holds the superposition of the possible products.
\State \Return $R$
\end{algorithmic}
\end{algorithm}

We can think of $C_j$ as a multiplication factor that is applied in each individual step, just as Regev would apply, for example, $b_1b_3b_4$.

We should mention that the classical postprocessing differs a bit, specifically the part we outlined in Chapter 10. With respect to complexity, given a circuit that implements something similar to $\textbf{Eq 13.2}$ with $S = n+ O(1)$ qubits and $G$ gates, the above algorithm can be executed with $S + (A/log\phi + 6 + O(1))n$ qubits and $O(\sqrt{n}G + n^{3/2})$ gates assuming an arbitrary multiplication circuit. Assuming we adapt the classical Van Der Hoeven algorithm for multiplication\cite{HarveyVanDerHoeven2021}, this is equivalent is equivalent to $O(nlogn)$ qubits and $O(n^{3/2} logn)$ gates, achieving a significant improvement over Regev's $O(n^{3/2})$ qubits.






